\abstract{РЕФЕРАТ}

Объем работы равен \formbytotal{lastpage}{страниц}{е}{ам}{ам}. Работа содержит \formbytotal{figurecnt}{иллюстраци}{ю}{и}{й}, \formbytotal{tablecnt}{таблиц}{у}{ы}{}, \arabic{bibcount} библиографических источников и \formbytotal{числоПлакатов}{лист}{}{а}{ов} графического материала. Количество приложений – 2. Графический материал представлен в приложении А. Фрагменты исходного кода представлены в приложении Б.

Перечень ключевых слов: онлайн-площадка, аренда вещей, модуль взаимодействия, чат, уведомления, генерация договоров, WebSocket, REST API, Django, Django REST Framework, Django Channels, PDF, push-уведомления, email-уведомления, брокер сообщений, Redis.

Объектом разработки является модуль взаимодействия с пользователем, включающий подсистемы чатов, уведомлений и автоматической генерации договоров аренды. Модуль обеспечивает обмен сообщениями в реальном времени, доставку уведомлений через WebSocket и электронную почту, а также формирование юридически значимых документов в формате PDF.

Целью выпускной квалификационной работы является разработка модуля взаимодействия с пользователями онлайн-площадки для сдачи вещей в аренду, включающего подсистему уведомлений, подсистему пользовательских чатов и подсистему автоматической генерации договоров аренды.

В процессе разработки модуля взаимодействия проведён анализ предметной области онлайн-аренды и существующих решений в области встроенных мессенджеров, систем уведомлений и генераторов договоров. Выделены основные сущности (сообщения, уведомления, договоры) и установлены связи между ними. Разработана архитектура модуля с использованием REST API для управления данными и WebSocket для обмена сообщениями в реальном времени. Реализована подсистема чатов с поддержкой текстовых сообщений и вложений, подсистема уведомлений с доставкой через WebSocket и email, а также подсистема автоматической генерации договоров аренды в формате PDF с демонстрационным механизмом подписания. Выполнено тестирование разработанных компонентов и их интеграция с серверной частью платформы.

\selectlanguage{english}
\abstract{ABSTRACT}

The volume of work is \formbytotal{lastpage}{page}{}{s}{s}. The work contains \formbytotal{figurecnt}{illustration}{}{s}{s}, \formbytotal{tablecnt}{table}{}{s}{s}, \arabic{bibcount} bibliographic sources and \formbytotal{числоПлакатов}{sheet}{}{s}{s} of graphic material. The number of applications is 2. The graphic material is presented in annex A. Source code fragments are presented in annex B.

The list of keywords: online marketplace, rental of things, user interaction module, chat, notifications, contract generation, WebSocket, REST API, Django, Django REST Framework, Django Channels, PostgreSQL, PDF, push notifications, email notifications, message broker, Redis.

The object of development is a user interaction module for an online rental platform, including chat, notification and automatic contract generation subsystems. The module provides real-time messaging, delivery of notifications via WebSocket and email, as well as the generation of legally significant documents in PDF format.

The purpose of the final qualifying work is to develop a user interaction module for an online platform for renting out things, including a notification subsystem, a user chat subsystem, and an automatic rental contract generation subsystem.

During the development of the interaction module, an analysis of the subject area of online rental and existing solutions in the field of built-in messengers, notification systems and contract generators was carried out. The main entities (chats, messages, notifications, contracts, deals) were identified and relationships between them were established. The module architecture was developed using REST API for data management and WebSocket (Django Channels) for real-time messaging. A chat subsystem with support for text messages and attachments, a notification subsystem with delivery via WebSocket and email, and an automatic rental contract generation subsystem in PDF format with a demonstration signing mechanism were implemented. Testing of the developed components and their integration with the server part of the platform was performed.

\selectlanguage{russian}